\documentclass[amsfonts,12pt]{article}

%\usepackage{amsbsy}
\usepackage{amssymb}
\usepackage{amsthm}
\usepackage{amsgen}
\usepackage{amsfonts}

\newtheorem{theorem}{Theorem} 
\newtheorem{corollary}[theorem]{Corollary} 
\newtheorem{conjecture}[theorem]{Conjecture} 
\newtheorem{lemma}[theorem]{Lemma} 
\newtheorem{proposition}[theorem]{Proposition} 
\newtheorem{definition}[theorem]{Definition} 
\newtheorem{example}[theorem]{Example} 
\newtheorem{axiom}{Axiom} 
\newtheorem{remark}{Remark} 
\def\fff{\mathbb{F}}
\def\aaa{\mathbb{A}}
\def\qqq{\mathbb{Q}}
\def\rrr{\mathbb{R}}
\def\ccc{\mathbb{C}}
\def\zzz{\mathbb{Z}}
\def\hhh{{\cal H}}

\begin{document}

\author{David Joyner}
\title{The representation theory of finite commutation 
transitive groups\\
PRELIMINARY VERSION}
\date{8-5-2000}
\maketitle

\begin{abstract}
We investigate certain induced representations of finite 
commutation transitive groups.
We investigate which irreducible characters of such a group
are generated by certain monomial representations.
The paper concludes with some remarks on the
profinite case.
\end{abstract}

\section{Introduction}

Let $G$ be a finite group and let $R$ be a reflexive,
symmetric relation on $G$. Roughly speaking,
we call $G$ {\bf $R$-transitive}
if $R$ is an equivalence relation on $G$. 

Let $R$ be the relation
given by ``$(x,y) \in R$ if and only if $x$ commutes with $y$''.
We call $G$ {\bf commutative transitive}
if $R$ is an equivalence relation on $G-\{1\}$. This notion
was first introduced in 1925 by Weisner \cite{W} but more recently
was (independently and from a different perspective) 
developed in \cite{FGRS}, \cite{Wu}. Such groups were classified by 
Weisner, Suzuki, Wall, and Wu. Any commutative transitive group
must either be abelian or have trivial center.

Let $Z=Z(G)$ denote the center of $G$ and let $R$ be the
commutivity relation above. If $R$ is an equivalence relation on 
the set $G-Z$ then
we call $G$ {\bf commutation transitive}. I do not know
the origin of this term, though some older references
may be found in \cite{LR}.

\begin{theorem} (\cite{S}, \cite{W}, \cite{Wu}) 
A finite solvable group $G$ is 
commutative transitive if and only if one of the following holds:
\begin{enumerate}
\item
$G$ is abelian,

\item
$G$ is non-abelian and contains both a normal
abelian subgroup $N$ and a cyclic group $H$ of
fixed-point-free autmorphisms of $N$ for which 
$G=H\ltimes N$.

\end{enumerate}

A finite simple group $G$ is 
commutative transitive if and only if one of the following holds:
\begin{enumerate}
\item
$G$ is cyclic,

\item
$G\cong PSL_2({\Bbb{F}}_{2^k})$, for some $k\geq 2$.
\end{enumerate}
\end{theorem}

\begin{remark} Weisner \cite{W} gives 
$G\cong PSL_2({\Bbb{F}}_{2^k})$, for some $k> 2$,
but $PSL_2({\Bbb{F}}_{2^2})$ is a simple group of order 60, 
hence must be isomorphic to $A_5$ which is
commutative transitive. 

Wu \cite{Wu} pointed out that
one of Weisner's proofs relied on a faulty result,
hence contained a gap. However, the result stated 
above still holds.
\end{remark}

\begin{example}
If $q$ is an odd prime power then $G=PGL(2,\fff_q)$ is
commutative transitive but not simple. (Recall
(1) $PSL(m,\fff_q)=PGL(m,\fff_q)$ if $gcd(m,q-1)=1$,
and (2) $PSL(2,\fff_q)$ is simple if and only if 
$q>3$ \cite{R}.).
\end{example}

\begin{theorem} (\cite{F}, \cite{LR}) 

\begin{enumerate}
\item
A free group $G$ is commutation transitive.
\item
$GL(2,k)$, where $k$ is any field, is commutation transitive.
\end{enumerate}
\end{theorem}

\begin{remark}
Let $G$ be commutation transitive and so

(a) $x\sim x$, for all $x\in G-Z$,

(b) $x_1\sim x_2$ implies $ x_2\sim x_1$, for all $x_1,x_2\in G-Z$,

(c) $x_1\sim x_2,\ x_2\sim x_3$ implies $x_1\sim x_3$, for all 
$x_1,x_2,x_3\in G-Z$.

\noindent
Now we extend this to $G/Z$. First,
define $\overline{x}=xZ$, for $x\in G$.
Note, for all $x,y\in G-Z$,
$xy=yx$ if and only if, for all $z_1,z_2\in Z$, we have
$xz_1$ commutes with $yz_2$. In other words,
\[
x\sim y \iff \forall z_1,z_2\in Z,\ xz_1\sim yz_2.
\]
Now define $\sim $ (with apologies for the abuse of
notation) on $(G-Z)/Z$ by
$\overline{x}\sim \overline{y}$ if and only if, 
for {\it some} $z_1,z_2\in Z$, we have
$xz_1$ commutes with $yz_2$.
This is well-defined since in fact we have shown
$x\sim y$ if and only if $\overline{x}\sim \overline{y}$.
Since $1\in G/Z$ may be identified with $Z$, 
this new $\sim$ is a relation on the 
subset $(G-Z)/Z\subset G/Z$.
It is clearly reflexive and symmetric. However, it
is possible for $\overline{x}_1$ to commute with 
$\overline{x}_2$ even
though $x_1, x_2$ do not commute
\footnote{Take 
$x_1=\left(\begin{array}{cc}-1 & 0\\ 0 & 1\end{array}
\right) $ and $x_2=\left(\begin{array}{cc}0 & 1\\ 1 & 0\end{array}
\right) $. Then $x_1^{x_2}=-x_1$ but
$\overline{x}_1^{\overline{x}_2}=\overline{x}_1$.},
so transitivity can fail.
\end{remark}


\section{Background}

 One can create a graph $\Gamma = (V,E)$ with vertices 
  the set of elements of $G-Z$ and whose edges are given by 
  the following rule: $v,v' \in V$ are connected by an edge
  if and only if they commute. 

\begin{lemma} Assume $G$ is commutation transitive. 

  (a) Each connected component of $\Gamma$ is a complete
  graph. 

  (b) The connected components of $\Gamma$ and the set of centralizers
$\{Cent(g,G)\ |\ g\in G-Z\}$ are in a natural one-to-one correspondence.
\end{lemma}

\begin{proof} Trivial. \end{proof}


    Let $g \in G$, let
\[
         Cent(g,G) = \{ x \in G\ |\ xg = gx\}
\]
denote the centralizer of $g$ in $G$, let
\[
         Conj(X,Y) = \{ y^{-1}xy\ |\ x\in X,\ y\in Y\},
\ \ \ \ X\subset G,\ Y\subset G,
\]
denote the conjugacy class of $X$ in $Y$, and let
\[
        N_G(H) = \{ x \in G\ |\ xH = Hx\}
\]
denote the normalizer of $H$ in $G$. If $X=\{x\}$ and 
$Y=\{y\}$ are singleton sets, we write $Conj(x,y)=y^{-1}xy$.
Let $G_*$ denote the set
of conjugacy classes in $G$. It is easy to show that
the conjugate of a centralizer is the centralizer of the 
conjugate:
\begin{equation}
         y^{-1}Cent(g,G)y = Cent(y^{-1}gy,G),
\label{eq:*}
\end{equation}
for all $y \in G$. Also, there is a natural bijection
between $Conj(g,G)$ and $G/Cent(g,G)$ given by 
\[
  x^{-1}gx  \longleftrightarrow xCent(g,G).
\]
  
\begin{lemma} Assume $G$ is commutation transitive. Let 
$g_0\in G-Z$. If $g_1\in Cent(g_0,G)-Z$ then 
$Cent(g_0,G)=Cent(g_1,G)$. In
particular, the centralizer of a centralizer is itself.
\label{lemma:cent}
\end{lemma}

\begin{proof} This is a consequence of the fact that commutativity
is an equivalence relation on $G-Z$.   
\end{proof}

In particular, any two centralizers of a commutative transitive group
are either disjoint groups or equal to each other.

Suppose that $C=Cent(g_0,G)$ is a proper subgroup
and that it is equal to its normalizer in $G$:
$C=N_G(C)$. By the previous lemma, such a group is 
a Frobenius subgroup. In this case, by Exercise 7.3 
in Serre \cite{Se}, $G$ must have
a normal subgroup of size at least $2$. In particular,
if $C=N_G(C)$ then $G$ cannot be simple.

%\begin{remark}
%It is known that if $H$ is any proper subgroup of $G$ then 
%$\cup_{x\in G} Conj(H,x)$ is a proper subset of $G$ 
%(\cite{R}, Exercise 3.13)
%\footnote{This follows from the previous statement if
%$G$ is finite. If
%$H=N_G(H)$ then a counting argument shows 
%$G=\cup_{x\in G} Conj(H,x)$ cannot occur.
%If $H\not= N_G(H)$ then a different counting argument shows 
%$|\cup_{x\in G} Conj(H,x)|<|G|$. If $G$ is infinite then
%the analogous result fails to hold (thanks to
%Dennis Spellman for providing a counterexample).}.
%
%It also is known that the number of
%distinct conjugates is equal to the index of $N_G(H)$
%in $G$ (\cite{R}, Theorem 3.3):
%\[
%|\{ Conj(H,x)\ |\ x\in G\}| = |G/N_G(H)|.
%\]
%\end{remark}


\section{Representations of the centralizers}

 Let $H^*$ denote the set of equivalence classes of 
irreducible representations of a finite group $H$.

\begin{lemma} $G$ is commutation transitive if and only if
each centralizer of a non-central element is abelian.
\end{lemma}

\begin{remark} This is actually the form in which Weisner
introduced the notion (of commutative transitive).
\end{remark}

\begin{proof} This is an immediate consequence of the definition.
\end{proof}

   This lemma, of course, 
implies that every irreducible representation of 
a centralizer is 1-dimensional. 
    
\section{Representations of $G$}

\subsection{The solvable case}

If $G$ is a solvable commutative transitive group whcih is not cyclic
then the methif of ``little groups'', due to 
Wigner and Mackey, applies to construct all
the irreducible representations of $G$.
See \S 8.2 of Serre \cite{Se}.

\subsection{The general case}

The representations of finite commutative transitive groups
were also studied by Suzuki \cite{S}. His description of 
$G^*$ was different from that given here.

We first state a criterion for irreducibility
of induced representations. (Note this criterion holds even in the 
profinite case.)

\begin{lemma} (Mackey's irreducibility criterion)
Let $G$ be a locally compact, totally disconnected, separable,
unimodular group. Let $Z$ be the center of $G$ and let $H$ be an open
subgroup of $G$ containing $Z$ which is compact mod $Z$. Let
$\pi$ be a smooth irreducible representation of $H$ satisfying
\begin{itemize}
\item[(1)]
for all $g\in G-H$, $\pi|_{H\cap H^g}$ and
$\pi^g|_{H\cap H^g}$ are disjoint,
\item[(2)]
the modulus $|\chi|$ of the central character $\chi$ of
$\pi$ extends to character of $G$.
\end{itemize}
Then $ind_H^G\pi$ is an irreducible admissible representation. 
\end{lemma}

This form goes back to Mauntner (and may be found in \cite{J}
for example).
 
\begin{remark}
(a) If $G$ is a finite group then condition
(2) is trivially true since $\chi$ is of finite order.

(b) Let $G=GL(2,\fff_q)$ or $G=PSL(2,\fff_q)$ 
and let $H$ be a maximal torus. Then
$H\cap H^g=H$ or $H\cap H^g=Z\cong \fff_q^\times$. In particular
condition (1) is always false
\footnote{Disjoint groups {\it do not} have disjoint representations!
}
.

\end{remark}

If $\pi\in G^*$, let ${\rm tr}\, \pi$ denote its character:
${\rm tr}\,\pi(g)={\rm tr}(\pi(g))$, for $g\in G$.
The Frobenius formula for the character of an induced representation
in this case becomes particularly simple:

\begin{lemma} Let $G$ be a finite commutation transitive group,
$C$ a centralizer in $G$, $\mu\in C^*$ a character,
and let $\pi=ind_C^G\mu$. Then 
\[
{\rm tr}\, \pi(g)=
\left\{
\begin{array}{cc}
\sum_{n\in N_G(C)}\mu(n^{-1}an), & a=a_g\in Conj(g,G)\cap C,\ g\not= 1\\
0, &Conj(g,G)\cap C =\emptyset,\ g\not= 1,\\
\frac{|G|}{|C|},& g=1.
\end{array}
\right.
\]
\end{lemma}

\begin{proof} This is a consequence of the formula
for induced characters \cite{Se}, chapter 7, and
Lemma \ref{lemma:cent}.
\end{proof}


\begin{corollary} (of Proposition)
Let $G=PGL(2,\fff_q)$, let $T$ be a maximal torus
of $G$ (more precisely, the $\fff_q$-valued points
of a maximal torus of $PGL(2)$, and let $H=N_G(T)$
denote its normalizer. If $\pi$ is a character of $T$ 
such that $\pi\not= \pi^h$ for all $h\in H-T$ then
$ind_T^G \pi$ is reducible.
\end{corollary}

\begin{proof}
The Frobenius formula and the hypothesis implies that
all the terms in 
\[
\begin{array}{c}
(ind_T^G \pi,ind_T^G \pi)=\frac{1}{|G|}\sum_{g\in G}
({\rm tr}ind_T^G \pi(g))({\rm tr}ind_T^G \pi(g^{-1})\\
=\frac{1}{|G|}\sum_{g\in G}\sum_{n_1\in N_G(T)}
\sum_{n_2\in N_G(T)}\mu(n_1^{-1}a_g n_1)\mu(n_2^{-1}a_g^{-1}n_2)
\end{array}
\]
cancel except for those for which the character
\[
a=a_g\longmapsto  \mu(n_1^{-1}a_g n_1)\mu(n_2^{-1}a_g^{-1}n_2)
\]
is trivial. The contribution to these terms is positive;
in fact, those where $g=1$ yield a contribution $>1$.
\end{proof}

\begin{definition}
For any subgroup $H$ of $G$, define the
{\bf support} of $H$ in $G_*$ to be
\[
supp(H)=\{\gamma\in G_*\ | \ \gamma\cap H\not= \emptyset\}.
\]
\end{definition}

Let $\fff_2^{G_*}=\fff_2[G_*]$ denote the vector space over $\fff_2$
formally spanned by the elements of $G_*$. This may be identified
in a natural way with the power set offfff $G_*$ (i.e., the set
of all subsets of $G_*$.) In particular, each subgroup $H$
of $G$ may be associated in a natural way a vector in $\fff_2^{G_*}$.
This yields a natural map
\begin{equation}
\tau:L(G)\rightarrow \fff_2^{G_*},
\end{equation}
where $L(G)$ is the subgroup lattice of $G$.

Let $G$ be a finite commutation transitive group.
For each $\gamma\in G_*$, pick $ g_0\in \gamma$,
compute $A = Cent(g_0,G)$, and let $\chi_A \in A^*$
be the trivial character.
Let ${\cal C}$ denote a set of centralizers $A$ which 
arise in this way and satisfy the following property:
\[
\{\tau(A)\ |\ A\in {\cal C}\}\subset \fff_2^{G_*}
\]
is a linearly independent set.

\begin{theorem} 
\label{thrm:main}
The set of characters of representations of $G$ obtained
by inducing one character from each centralizer in ${\cal C}$,
\[
\{{\rm tr} ind_A^G(\chi)\ |\ \chi=\chi_\gamma \in A^*\ , A\in {\cal C}\},
\]
is linearly independent.
\end{theorem}

\begin{proof}
In fact, the functions are supported on pairwise
distinct sets
\footnote{
This is a consequence of Frobenius' formula for characters
of induced representations.}, so it is impossible for any linear
combination of them to add to zero.

\end{proof}

For examples in the case $G=A_5$, see \cite{J2}.
It is possible for distinct centralizers to have 
identical support.

The following guess is based on computer computations for 
$G=PSL(2,2^k)$,
$k=2,3,4,5$ using GAP. It seems to place a restriction on the
number of monomial representations of $G$.

\begin{conjecture}
If $G$ is a finite simple commutative transitive group then
the span of the vectors $\tau(Cent(G,g))$, $g\in \gamma\in G_*$,
is bounded independent of $|G|$. In fact,
it appears that 
\[
{\rm dim\ Span}\ \{\tau(Cent(G,g))\ |\  g\in \gamma\in G_*\}
\leq 4.
\]
\end{conjecture}

The above theorem says that, in principle at least, all the 
representations in $G^*$ can be obtained from linear combinations of 
monomial representations. This type of result goes back to 
Gelfand and Graev.

\section{Profinite groups}

We conclude this papers with some remarks.

There is an analog of 
Theorem \ref{thrm:main} for profinite commutative 
transitive groups.

{\bf Acknowledgements}:
I thank Tony Gaglione, Dennis Spellman, and Jill McGowan 
for helpful conversations.

\begin{thebibliography}{99}  

\bibitem[B]{B} A. Borel, R. Carter, C. Curtis, N, Iwahori, 
T. Springer, and R. Steinberg, {\bf Seminar on algebraic
groups and related finite groups},  Springer
Lecture Notes in Mathematics, vol 131, Springer-Verlag,
2nd printing 1986

\bibitem[F]{F} B. Fine, ``On power conjugacy and $SQ$-universality 
for Fuchsian groups and Klienian groups,'' preprint 1983

\bibitem[FGRS]{FGRS} B. Fine, A. Gaglione,
G. Rosenberger, D. Spellman, ``The commutative transitive kernel'',
to appear in Algebra Colloquium

\bibitem[J1]{J} D. Joyner, ``A correspondence for the supercuspidal
representations of $\overline{SL_2(F)}$, $F$ p-adic,'' 
Archiv der Math. \underbar{273}(1999)1-9

\bibitem[J2]{J2} -----, ``Representation theory of $A_5$ using GAP,''
preprint, 2000

\bibitem[LR]{LR} F. Levin, G. Rosenberger, ``On power-commutative
and commutative-transitive groups,'' in
{\bf Proc. of groups - St. Andrews, 1985},
London Math. Soc. Lecture Notes, 121, Cambridge Univ.
Press, 1986.

\bibitem[R]{R} J. J. Rotman, {\bf An introduction to the
theory of groups}, 4th ed, Springer-Verlag, Grad Texts in Math 
\underbar{148}, 1995

\bibitem[Se]{Se} J.-P. Serre, 
{\bf Linear representations of finite groups},
Springer-Verlag, 1977

\bibitem[S]{S} M. Suzuki, ``The nonexistence of a certain type of
simple groups of even order'', Proc A.M.S. \underline{8}(1957)
686-695

\bibitem[W]{W} L. Weisner, ``Groups in which the normalizer of 
every element except the identity is abelian'', Bull. A.M.S. 
\underline{31}(1925)413-416

\bibitem[Wu]{Wu} Y.-F. Wu, ``Groups in which commutivity is a
transitive relation,'' J. of Alg. \underline{207}(1998)165-181
\end{thebibliography}
\end{document}
b